%================================ ABSTRACT ====================================
\clearpage
\thispagestyle{plain}
\addcontentsline{toc}{chapter}{Abstract}
\begin{center}
{\Large\bfseries ABSTRACT\par}
\end{center}
\vspace{4mm}

A learner who wants a quick visual answer to ``how does this work'' usually has
to stitch together text from one source, a static picture from another, and a
separate video from a third. This thesis presents \textbf{TutorAI}, a complete,
working system that turns a single typed topic into a short lesson that
explains itself. Each lesson consists of three tightly coupled parts: a vector
diagram (SVG) whose meaningful parts carry stable, semantic identifiers; a
narration script split into ordered segments, where every segment names the
diagram parts it talks about; and one synthesized speech clip per segment with
a precisely measured duration. On an Android device, playback highlights the
referenced diagram parts while the matching audio plays, recreating the
``point and explain'' behaviour of a human tutor. Once downloaded, a lesson
replays fully offline and is kept in an on-device history.

The central engineering challenge is reliability: a large language model will
happily write narration that references a diagram part that does not exist,
which silently breaks the highlighting. TutorAI solves this with an agentic
generation pipeline built on LangGraph that separates planning, drawing,
critique-and-refine of the drawing, narration, and validation into distinct
nodes, and places a deterministic, non-LLM validator as a hard gate in front
of the cache: a lesson is never stored unless its SVG is well formed and every
referenced identifier really exists. Bounded refine and repair loops raise
quality while capping cost and latency. The backend is a stateless FastAPI
service with a shared, topic-keyed lesson cache, so a popular topic is
generated once and then served instantly. Both the LLM (Gemini) and the
text-to-speech engine sit behind swappable provider adapters, making model
choice a configuration concern. The Android client follows Clean Architecture
with MVVM, uses Room and a private file store for offline persistence, and
drives segment-synchronized highlighting through a WebView-hosted SVG and
ExoPlayer playlist.

The realized system was verified end to end: live generation produced a valid,
narrated, highlight-consistent lesson; the validation gate was shown to reject
and repair broken lessons; a hermetic suite of 48 key-free backend tests runs
in continuous integration; and the containerized backend and installable
Android build are produced by automated pipelines. The thesis describes the
methodology, architecture, implementation, and verified results, and discusses
limitations and future work including formal learning-outcome studies.

\vspace{6mm}
\noindent\textbf{Keywords:} generative AI, large language models, LLM agents,
intelligent tutoring systems, multimedia learning, SVG generation,
text-to-speech, FastAPI, LangGraph, Android, offline-first, software
architecture.
